
\subsection{Corrections des Exercices}

\begin{correctionbox}[Correction Exercice 1 : Définition (Paramètres)]
1.  Par définition, si $Y \sim \text{Log-}\mathcal{N}(\mu, \sigma^2)$, alors $X = \ln(Y)$ suit une loi normale $\mathcal{N}(\mu, \sigma^2)$.
    Ici, $X \sim \mathcal{N}(3, 16)$.
2.  $E[X] = \mu = 3$.
    $\text{Var}(X) = \sigma^2 = 16$.
\end{correctionbox}

\begin{correctionbox}[Correction Exercice 2 : Définition (Inverse)]
Si $X \sim \mathcal{N}(0, 1)$, alors $Y = e^X$ suit une loi log-normale avec $\mu=0$ et $\sigma^2=1$.
La notation est $Y \sim \text{Log-}\mathcal{N}(0, 1)$.
\end{correctionbox}

\begin{correctionbox}[Correction Exercice 3 : Intuition (Somme vs Produit)]
La loi normale résulte de l'addition de nombreux petits effets (TCL).
La loi log-normale résulte de la multiplication de nombreux petits facteurs.
La croissance d'une population de bactéries est un processus multiplicatif ($P(t) = P(t-1) \times F_t$).
En prenant le logarithme, $\ln(P(t)) = \ln(P(0)) + \sum \ln(F_i)$, on obtient une somme qui tend vers une loi normale. Donc $P(t)$ lui-même tend vers une loi log-normale.
De plus, $P(t)$ doit être positif, ce que garantit la loi log-normale (support $(0, \infty)$).
\end{correctionbox}

\begin{correctionbox}[Correction Exercice 4 : Intuition (Transformation)]
La transformation $Y = e^X$ n'est pas linéaire. Elle est convexe.
Elle "étire" la partie droite de la distribution normale ($x>0 \implies e^x$ croît exponentiellement) et "compresse" la partie gauche ($x<0 \implies e^x$ s'écrase vers 0).
Une distribution symétrique (Normale) transformée par une fonction asymétrique (Exponentielle) donne une distribution asymétrique (Log-Normale).
\end{correctionbox}

\begin{correctionbox}[Correction Exercice 5 : PDF (Propriétés)]
Le support de la loi log-normale $Y$ est $(0, \infty)$. La variable ne peut jamais être négative ou nulle.
Par conséquent, $P(Y \le 0) = 0$.
\end{correctionbox}

\begin{correctionbox}[Correction Exercice 6 : Espérance]
$Y \sim \text{Log-}\mathcal{N}(\mu=4, \sigma^2=2)$.
$E[Y] = e^{\mu + \sigma^2/2} = e^{4 + 2/2} = e^{4 + 1} = e^5$.
\end{correctionbox}

\begin{correctionbox}[Correction Exercice 7 : Médiane]
$Y \sim \text{Log-}\mathcal{N}(\mu=4, \sigma^2=2)$.
$\text{Med}(Y) = e^{\mu} = e^4$.
\end{correctionbox}

\begin{correctionbox}[Correction Exercice 8 : Mode]
$Y \sim \text{Log-}\mathcal{N}(\mu=4, \sigma^2=2)$.
$\text{Mode}(Y) = e^{\mu - \sigma^2} = e^{4 - 2} = e^2$.
\end{correctionbox}

\begin{correctionbox}[Correction Exercice 9 : Relation d'Ordre]
Les valeurs sont :
$\text{Mode} = e^2 \approx 7.39$
$\text{Médiane} = e^4 \approx 54.60$
$\text{Espérance} = e^5 \approx 148.41$
On vérifie bien $\text{Mode} < \text{Médiane} < \text{Espérance}$, ce qui est la signature d'une asymétrie à droite.
\end{correctionbox}

\begin{correctionbox}[Correction Exercice 10 : Variance]
$Y \sim \text{Log-}\mathcal{N}(\mu=1, \sigma^2=1)$.
$\text{Var}(Y) = (e^{\sigma^2} - 1) \cdot e^{2\mu + \sigma^2}$
$\text{Var}(Y) = (e^{1} - 1) \cdot e^{2(1) + 1} = (e - 1)e^3$.
\end{correctionbox}

\begin{correctionbox}[Correction Exercice 11 : Moments (Problème Inverse)]
1.  $\text{Med}(Y) = e^{\mu}$. On nous donne $\text{Med}(Y) = e^5$. Donc $\mu = 5$.
2.  $E[Y] = e^{\mu + \sigma^2/2}$. On nous donne $E[Y] = e^{5.5}$.
    Donc, $e^{5.5} = e^{\mu + \sigma^2/2} = e^{5 + \sigma^2/2}$.
    En égalant les exposants : $5.5 = 5 + \sigma^2/2$.
    $0.5 = \sigma^2/2 \implies \sigma^2 = 1$.
\end{correctionbox}

\begin{correctionbox}[Correction Exercice 12 : Impact de $\sigma$ sur l'Espérance]
$E[A] = e^{\mu_A + \sigma_A^2/2} = e^{0 + 1/2} = e^{0.5}$.
$E[B] = e^{\mu_B + \sigma_B^2/2} = e^{0 + 2/2} = e^{1}$.
$E[B] > E[A]$. L'espérance $E[Y] = e^{\mu + \sigma^2/2}$ croît avec $\sigma^2$. Cela est dû à l'inégalité de Jensen ($\mathbb{E}[e^X] \geq e^{\mathbb{E}[X]}$) : une plus grande variance (incertitude) sur $X=\ln(Y)$ augmente l'espérance de $Y=e^X$ à cause de la convexité de l'exponentielle.
\end{correctionbox}

\begin{correctionbox}[Correction Exercice 13 : Calcul de CDF (Simple)]
On cherche $P(Y \le e^5)$ pour $Y \sim \text{Log-}\mathcal{N}(5, 1)$.
$$P(Y \le e^5) = P(\ln(Y) \le \ln(e^5)) = P(X \le 5)$$
Puisque $X = \ln(Y) \sim \mathcal{N}(\mu=5, \sigma^2=1)$, on cherche la probabilité que $X$ soit inférieure à sa propre moyenne. Par symétrie de la loi normale, c'est 0.5.
Formellement : $Z = \frac{5 - 5}{1} = 0$. $P(Z \le 0) = \Phi(0) = 0.5$.
\end{correctionbox}

\begin{correctionbox}[Correction Exercice 14 : Calcul de CDF]
$Y \sim \text{Log-}\mathcal{N}(\mu=3, \sigma=2)$.
$P(Y \le e^4) = P(\ln(Y) \le \ln(e^4)) = P(X \le 4)$, où $X \sim \mathcal{N}(3, 4)$.
Standardisation : $Z = \frac{X - \mu}{\sigma} = \frac{4 - 3}{2} = \frac{1}{2} = 0.5$.
$P(X \le 4) = P(Z \le 0.5) = \Phi(0.5) \approx 0.6915$.
\end{correctionbox}

\begin{correctionbox}[Correction Exercice 15 : Calcul de Probabilité (Queue Droite)]
$Y \sim \text{Log-}\mathcal{N}(0, 1)$.
$P(Y > 1) = 1 - P(Y \le 1) = 1 - P(\ln(Y) \le \ln(1)) = 1 - P(X \le 0)$.
$X \sim \mathcal{N}(0, 1)$. $P(X \le 0) = \Phi(0) = 0.5$.
$P(Y > 1) = 1 - 0.5 = 0.5$. (Logique : $e^0=1$ est la médiane, donc 50% est au-dessus).
\end{correctionbox}

\begin{correctionbox}[Correction Exercice 16 : Calcul de Probabilité (Queue Gauche)]
$Y \sim \text{Log-}\mathcal{N}(1, 1)$.
$P(Y \le e^{-1}) = P(\ln(Y) \le \ln(e^{-1})) = P(X \le -1)$, où $X \sim \mathcal{N}(1, 1)$.
Standardisation : $Z = \frac{X - \mu}{\sigma} = \frac{-1 - 1}{1} = -2$.
$P(X \le -1) = P(Z \le -2) = \Phi(-2) = 1 - \Phi(2) \approx 1 - 0.9772 = 0.0228$.
\end{correctionbox}

\begin{correctionbox}[Correction Exercice 17 : Calcul de Probabilité (Intervalle)]
$Y \sim \text{Log-}\mathcal{N}(10, 9)$. ($X \sim \mathcal{N}(10, 9)$, $\sigma=3$).
$P(e \le Y \le e^{16}) = P(\ln(e) \le \ln(Y) \le \ln(e^{16})) = P(1 \le X \le 16)$.
Standardisation :
$z_1 = \frac{1 - 10}{3} = -3$.
$z_2 = \frac{16 - 10}{3} = 2$.
$P(1 \le X \le 16) = P(-3 \le Z \le 2) = \Phi(2) - \Phi(-3)$.
$\Phi(-3) = 1 - \Phi(3) \approx 1 - 0.9987 = 0.0013$.
$P = \Phi(2) - \Phi(-3) \approx 0.9772 - 0.0013 = 0.9759$.
\end{correctionbox}

\begin{correctionbox}[Correction Exercice 18 : Application (Revenu)]
$Y \sim \text{Log-}\mathcal{N}(3, 1)$. ($X \sim \mathcal{N}(3, 1)$, $\sigma=1$).
$P(Y < e^2) = P(\ln(Y) < \ln(e^2)) = P(X < 2)$.
Standardisation : $Z = \frac{2 - 3}{1} = -1$.
$P(X < 2) = P(Z < -1) = \Phi(-1) = 1 - \Phi(1) \approx 1 - 0.8413 = 0.1587$.
\end{correctionbox}

\begin{correctionbox}[Correction Exercice 19 : Problème Inverse (Médiane)]
On cherche $y$ tel que $P(Y \le y) = 0.5$.
C'est la définition de la médiane. $\text{Med}(Y) = e^{\mu} = e^5$.
Donc $y = e^5$.
\end{correctionbox}

\begin{correctionbox}[Correction Exercice 20 : Problème Inverse (Percentile)]
$Y \sim \text{Log-}\mathcal{N}(0, 1)$. On cherche $y$ tel que $P(Y \le y) = 0.8413$.
$P(\ln(Y) \le \ln(y)) = 0.8413$.
$P(X \le \ln(y)) = 0.8413$, où $X \sim \mathcal{N}(0, 1)$.
$P(Z \le \ln(y)) = 0.8413$.
On cherche $z = \ln(y)$ tel que $\Phi(z) = 0.8413$. D'après la table, $z=1$.
$\ln(y) = 1 \implies y = e^1 = e$.
\end{correctionbox}

\begin{correctionbox}[Correction Exercice 21 : Problème Inverse (Percentile)]
$Y \sim \text{Log-}\mathcal{N}(2, 9)$. ($X \sim \mathcal{N}(2, 9)$, $\sigma=3$).
On cherche $y$ tel que $P(Y \le y) \approx 0.9772$.
$P(\ln(Y) \le \ln(y)) = P(X \le \ln(y)) = 0.9772$.
Standardisation : $P\left(Z \le \frac{\ln(y) - 2}{3}\right) = 0.9772$.
On cherche $z$ tel que $\Phi(z) = 0.9772$. D'après la table, $z=2$.
$\frac{\ln(y) - 2}{3} = 2 \implies \ln(y) - 2 = 6 \implies \ln(y) = 8 \implies y = e^8$.
\end{correctionbox}

\begin{correctionbox}[Correction Exercice 22 : Problème Inverse (Quantile bas)]
$Y \sim \text{Log-}\mathcal{N}(10, 4)$. ($X \sim \mathcal{N}(10, 4)$, $\sigma=2$).
On cherche $y$ tel que $P(Y > y) \approx 0.9938$.
$P(Y \le y) = 1 - 0.9938 = 0.0062$.
$P(X \le \ln(y)) = 0.0062$.
Standardisation : $P\left(Z \le \frac{\ln(y) - 10}{2}\right) = 0.0062$.
On cherche $z$ tel que $\Phi(z) = 0.0062$.
La table ne donne pas cette valeur, mais $\Phi(2.5) \approx 0.9938$.
Par symétrie, $\Phi(-2.5) = 1 - \Phi(2.5) \approx 1 - 0.9938 = 0.0062$.
Donc, $z = -2.5$.
$\frac{\ln(y) - 10}{2} = -2.5 \implies \ln(y) - 10 = -5 \implies \ln(y) = 5 \implies y = e^5$.
\end{correctionbox}

\begin{correctionbox}[Correction Exercice 23 : Modèle Financier (Paramètres)]
$X$ est une somme de 16 v.a. i.i.d. $x_i \sim \mathcal{N}(0.01, 0.04)$.
$X$ suit une loi normale.
$E[X] = \sum E[x_i] = t \cdot \mu = 16 \times 0.01 = 0.16$.
$\text{Var}(X) = \sum \text{Var}(x_i) = t \cdot \sigma^2 = 16 \times 0.04 = 0.64$.
Donc, $X \sim \mathcal{N}(0.16, 0.64)$.
\end{correctionbox}

\begin{correctionbox}[Correction Exercice 24 : Modèle Financier (Loi du Prix)]
$Y = S(16)/S(0)$. $X = \ln(Y) = \ln(S(16)/S(0))$.
D'après l'exercice 23, $X \sim \mathcal{N}(0.16, 0.64)$.
Par définition, $Y = e^X$ suit une loi log-normale.
$Y \sim \text{Log-}\mathcal{N}(\mu_X=0.16, \sigma_X^2=0.64)$.
\end{correctionbox}

\begin{correctionbox}[Correction Exercice 25 : Modèle Financier (Calcul de Prob.)]
$X = \ln(S(1)/S(0)) \sim \mathcal{N}(0.05, 0.09)$. (donc $\mu=0.05, \sigma=0.3$).
On cherche $P(S(1) < 100)$.
$P(S(1) < 100) = P(S(1)/S(0) < 100/100) = P(S(1)/S(0) < 1)$.
On applique le log :
$P(\ln(S(1)/S(0)) < \ln(1)) = P(X < 0)$.
On standardise $X$ :
$Z = \frac{X - \mu}{\sigma} = \frac{0 - 0.05}{0.3} = -0.05 / 0.3 \approx -0.167$.
$P(X < 0) = P(Z < -0.167)$.
C'est $\Phi(-0.167) = 1 - \Phi(0.167)$. Puisque $\Phi(0) = 0.5$, $\Phi(0.167)$ est légèrement supérieur à 0.5, donc $\Phi(-0.167)$ est légèrement inférieur à 0.5. La probabilité est légèrement inférieure à 50\%.
\end{correctionbox}
